\chapter{Podstawy teorii stabilności}
\label{ch:mod-stability}
\section{Motywacja i liczba typów}

Twierdzenie Ryll-Nardzewskiego
\ref{thm:mod-ryll-nardzewski} dotyczy skończoności
przestrzeni typów bez parametrów. Teoria stabilności
bada natomiast wzrost liczby typów w zależności
od mocy zbioru parametrów. W tym rozdziale
$T$ jest spełnialną teorią zupełną z modelami
nieskończonymi, a $S_n(A)$ ma znaczenie
z definicji \ref{def:mod-type-over-parameters}.

\begin{lemat}[Ogólne ograniczenie liczby typów]
\label{lem:mod-stability-type-bound}
Jeżeli $\mathfrak M\models T$, $A\subseteq M$
i $n\geq1$, to
\[
|S_n(A)|\leq 2^{\max(\aleph_0,|\mathcal L|,|A|)}.
\]
\end{lemat}
\begin{proof}
Dla $\mu=\max(\aleph_0,|\mathcal L|,|A|)$ język
z nazwami parametrów ma alfabet mocy co najwyżej
$\mu$. Skończonych ciągów symboli jest co najwyżej
$\mu$ (por. lemat \ref{lem:ls-formula-count}).
Typ jest podzbiorem zbioru formuł, zatem liczba
typów nie przekracza $2^\mu$.
\end{proof}

\begin{definicja}[Stabilność w mocy]
\label{def:mod-kappa-stability}
Dla $\kappa\geq\max(\aleph_0,|\mathcal L|)$ teorię
$T$ nazywamy $\kappa$-stabilną, jeżeli dla każdego
modelu $\mathfrak M\models T$ i każdego $A\subseteq M$
o $|A|\leq\kappa$ zachodzi $|S_1(A)|\leq\kappa$.
Teoria jest stabilna, jeśli jest $\kappa$-stabilna
dla pewnej takiej mocy $\kappa$.
W języku przeliczalnym $\aleph_0$-stabilność
nazywamy również $\omega$-stabilnością.
\end{definicja}

Liczymy wszystkie typy nad $A$, a nie tylko typy
realizowane w danym modelu. Typ może wymagać
rozszerzenia elementarnego, zgodnie z twierdzeniem
\ref{thm:mod-type-elementary-extension}.

\begin{twierdzenie}[Stabilność a typy krotek]
\label{thm:mod-stability-finite-tuples}
Jeżeli $T$ jest $\kappa$-stabilna, to dla każdego
$A$ o $|A|\leq\kappa$ i każdego $n\geq1$
zachodzi $|S_n(A)|\leq\kappa$.
\end{twierdzenie}
\begin{proof}
Stosujemy indukcję po $n$. Przypadek $n=1$
wynika z definicji. Ustalmy $q\in S_n(A)$
i wybierzmy jego realizację $\bar b$ w
rozszerzeniu elementarnym modelu zawierającego $A$.
Dla każdego typu $r\in S_{n+1}(A)$ ograniczającego
się do $q$ rozpatrzmy zbiór formuł
\[
r_{\bar b}(y)=
\{\varphi(\bar b,y):\varphi(\bar x,y)\in r\}.
\]
Parametry z $A$ pozostają bez zmian. Jeżeli
$\varphi_1,\ldots,\varphi_m\in r$, to formuła
$\exists y\,\bigwedge_{i=1}^{m}\varphi_i(\bar x,y)$
należy do $q$, ponieważ zachodzi przy pierwszych
$n$ współrzędnych każdej realizacji $r$.
Zatem jej podstawienie przy $\bar b$ jest prawdziwe.
Każdy skończony fragment $r_{\bar b}$ jest realizowany
w modelu zawierającym $\bar b$. Na mocy twierdzenia
\ref{thm:mod-type-elementary-extension} otrzymujemy
typ częściowy nad $B=A\cup\{b_1,\ldots,b_n\}$.
Każde jego zupełne rozszerzenie $s\in S_1(B)$
wyznacza $r$, gdyż $r$ rozstrzyga każdą formułę,
a $s$ zawiera odpowiednie podstawienie formuły
lub jej negacji.

Liczba typów ograniczających się do $q$ jest
więc nie większa niż $|S_1(B)|\leq\kappa$.
Z założenia indukcyjnego liczba możliwych $q$
jest nie większa niż $\kappa$.
Ponieważ $\kappa$ jest nieskończona,
$\kappa\cdot\kappa=\kappa$, co kończy dowód.
\end{proof}

\begin{przyklad}[Nieskończony zbiór z równością]
\label{ex:mod-stability-equality}
Niech $T$ będzie teorią nieskończonych zbiorów
w języku bez symboli pozalogicznych.
Z eliminacji kwantyfikatorów
(twierdzenie \ref{thm:mod-qe-equality})
typ w $S_1(A)$ jest określony przez informacje
o równości z nazwami elementów $A$.
Dla każdego $a\in A$ mamy dokładnie jeden typ
zawierający $x=a$. Istnieje też dokładnie jeden
typ zawierający wszystkie formuły $x\ne a$
dla $a\in A$: ich skończone fragmenty są
realizowalne, a eliminacja kwantyfikatorów
zapewnia jednoznaczność zupełnego rozszerzenia.
Są to wszystkie możliwości. Otrzymujemy
$|S_1(A)|=|A|+1$ w sensie arytmetyki mocy.
Teoria jest $\kappa$-stabilna dla każdej
nieskończonej mocy $\kappa$.
\end{przyklad}

\section{Własność porządku}

\begin{definicja}[Własność porządku]
\label{def:mod-order-property}
Formuła $\varphi(\bar x,\bar y)$ ma własność
porządku w teorii $T$, jeśli dla każdego $m\geq1$
istnieje model $\mathfrak M\models T$ i rodziny
krotek $(\bar a_i)_{i<m}$, $(\bar b_j)_{j<m}$,
takie że
\[
\mathfrak M\models\varphi(\bar a_i,\bar b_j)
\quad\Longleftrightarrow\quad i<j.
\]
Długości obu krotek zmiennych są dodatnie,
ustalone i niezależne od $m$.
\end{definicja}

Nie wymagamy, aby formuła definiowała porządek
na całym uniwersum. Wystarczy kodowanie dowolnie
długich skończonych porządków.

\begin{lemat}[Dowolny porządek indeksów]
\label{lem:mod-op-linear-index}
Jeżeli $\varphi$ ma własność porządku, to dla
każdego zbioru liniowo uporządkowanego $(I,<)$
istnieje model $T$ i rodziny krotek indeksowane
przez $I$, dla których
\[
\mathfrak M\models\varphi(\bar a_i,\bar b_j)
\quad\Longleftrightarrow\quad i<j.
\]
\end{lemat}
\begin{proof}
Dodajemy nowe stałe dla współrzędnych obu rodzin.
Do $T$ dołączamy zdania
$\varphi(\bar c_i,\bar d_j)$ dla $i<j$
oraz $\neg\varphi(\bar c_i,\bar d_j)$ dla $i\geq j$.
Skończony fragment tych wymagań używa skończenie
wielu indeksów. Porządkujemy je rosnąco i stosujemy
definicję \ref{def:mod-order-property}.
Z twierdzenia o zwartości \ref{thm:compactness}
istnieje model wszystkich wymagań.
Bierzemy jego redukt do języka pierwotnego.
\end{proof}

\begin{twierdzenie}[Własność porządku a $\omega$-stabilność]
\label{thm:mod-op-not-omega-stable}
Jeżeli język jest przeliczalny i pewna formuła
ma własność porządku w $T$, to $T$ nie jest
$\omega$-stabilna.
\end{twierdzenie}
\begin{proof}
Stosujemy lemat \ref{lem:mod-op-linear-index}
do $(\mathbb Q,<)$. Niech $A$ składa się ze
wszystkich współrzędnych krotek $\bar b_q$
dla $q\in\mathbb Q$. Zbiór $A$ jest przeliczalny.
Dla każdego $t\in\mathbb R$ tworzymy zbiór formuł
\[
\begin{aligned}
p_t={}&\{\varphi(\bar x,\bar b_q):q>t\}\\
 &{}\cup\{\neg\varphi(\bar x,\bar b_q):q\leq t\}.
\end{aligned}
\]
Każdy skończony fragment realizuje $\bar a_r$
dla odpowiedniego $r\in\mathbb Q$: wystarczy,
aby $r$ było nie mniejsze od wszystkich
indeksów ujemnych wymagań i mniejsze od
wszystkich indeksów dodatnich wymagań.
Jeśli oba rodzaje występują, największy
indeks ujemny jest mniejszy od najmniejszego
dodatniego; w przeciwnym razie korzystamy
z braku końców porządku $\mathbb Q$.
Z twierdzenia \ref{thm:mod-type-elementary-extension}
każdy $p_t$ jest typem częściowym nad $A$.
Wybieramy jego zupełne rozszerzenie $s_t$.
Dla $t<u$ wybieramy $q\in\mathbb Q$ z $t<q<u$.
Typy $s_t,s_u$ rozstrzygają formułę
$\varphi(\bar x,\bar b_q)$ przeciwnie.
Otrzymujemy co najmniej $2^{\aleph_0}$ typów
zupełnych w $|\bar x|$ zmiennych nad $A$.
Przeczy to $\omega$-stabilności na mocy
twierdzenia \ref{thm:mod-stability-finite-tuples}.
\end{proof}

\begin{przyklad}[Kategoryczność bez $\omega$-stabilności]
\label{ex:mod-dlo-not-omega-stable}
Teoria gęstych porządków liniowych bez końców
jest $\aleph_0$-kategoryczna (przykład
\ref{ex:mod-rn-dlo}), lecz formuła $x<y$
ma własność porządku. W $(\mathbb Q,<)$
wystarczy wziąć $a_i=b_i=i$ dla $0\leq i<m$.
Z twierdzenia \ref{thm:mod-op-not-omega-stable}
teoria nie jest $\omega$-stabilna.
\end{przyklad}


\section{Własność porządku w dowolnej mocy}

\begin{lemat}[Porządek o wielu przekrojach]
\label{lem:mod-op-many-cuts}
Dla każdej mocy nieskończonej $\kappa$ istnieje
zbiór liniowo uporządkowany $I$ oraz jego podzbiór
$J$ o mocy co najwyżej $\kappa$ i rodzina $H\subseteq I$
o mocy większej od $\kappa$, taka że każde dwa
różne elementy $H$ są rozdzielone elementem $J$.
\end{lemat}
\begin{proof}
Niech $\lambda$ będzie najmniejszą mocą, dla której
$2^\lambda>\kappa$. Istnieje ona z twierdzenia
Cantora i spełnia $\aleph_0\leq\lambda\leq\kappa$.
Moc traktujemy jako początkową liczbę porządkową.
Dla każdego $\alpha<\lambda$ mamy
$2^{|\alpha|}\leq\kappa$ z minimalności $\lambda$.

Rozważamy funkcje z $\lambda$ do $\{0,1,2\}$,
uporządkowane leksykograficznie: o kolejności
dwóch różnych funkcji decyduje ich wartość
na najmniejszym indeksie, na którym się różnią.
Jest to porządek liniowy; dla przechodniości
wystarczy porównać najmniejsze indeksy różnicy
dla dwóch kolejnych nierówności i rozpatrzyć
mniejszy z nich (lub ich wspólną wartość).

Niech $H=\{0,2\}^\lambda$. Dla każdego
$s\in\{0,2\}^\alpha$, $\alpha<\lambda$,
tworzymy funkcję $j_s$ równą $s$ przed $\alpha$
i równą $1$ od $\alpha$ do końca.
Niech $J$ będzie zbiorem wszystkich $j_s$,
a $I=J\cup H$ z odziedziczonym porządkiem.
Wówczas
\[
|J|\leq\sum_{\alpha<\lambda}2^{|\alpha|}
\leq\lambda\cdot\kappa=\kappa,
\qquad |H|=2^\lambda>\kappa.
\]
Jeżeli $h<h'$ należą do $H$, to na ich
pierwszym indeksie różnicy $\alpha$ mamy
$h(\alpha)=0$ oraz $h'(\alpha)=2$.
Ich wspólny fragment $s$ przed $\alpha$
wyznacza element $j_s$ spełniający
$h<j_s<h'$. To kończy dowód.
\end{proof}

\begin{twierdzenie}[Własność porządku wyklucza stabilność]
\label{thm:mod-op-unstable}
Jeżeli jakaś formuła ma własność porządku w $T$,
to $T$ nie jest stabilna w żadnej mocy
$\kappa\geq\max(\aleph_0,|\mathcal L|)$.
\end{twierdzenie}
\begin{proof}
Ustalmy taką moc $\kappa$ i wybierzmy $I,J,H$
z lematu \ref{lem:mod-op-many-cuts}.
Lemat \ref{lem:mod-op-linear-index} daje model
$\mathfrak M$ z krotkami $\bar a_i,\bar b_j$
dla $i,j\in I$, realizującymi wzorzec porządku
dla formuły $\varphi(\bar x,\bar y)$.
Niech $A$ będzie zbiorem współrzędnych krotek
$\bar b_j$ dla $j\in J$. Mamy $|A|\leq\kappa$,
ponieważ długość krotki $\bar y$ jest skończona.
Dla $h\in H$ bierzemy typ
$p_h=\operatorname{tp}^{\mathfrak M}(\bar a_h/A)$.
Jeśli $h<h'$, to istnieje $j\in J$ taki, że
$h<j<h'$. Formuła $\varphi(\bar x,\bar b_j)$
należy do $p_h$, a jej negacja do $p_{h'}$.
Zatem $p_h\ne p_{h'}$ i
\[
|S_{|\bar x|}(A)|\geq |H|>\kappa.
\]
Twierdzenie \ref{thm:mod-stability-finite-tuples}
wyklucza $\kappa$-stabilność.
\end{proof}

\section{Definiowalność typów nad modelem}

\begin{definicja}[Definiowalność typu względem formuły]
\label{def:mod-definable-type}
Niech $\mathfrak M\models T$,
$p\in S_n(M)$ oraz $\varphi(\bar x,\bar y)$
będzie formułą języka $\mathcal L$,
gdzie $|\bar x|=n$.
Typ $p$ jest definiowalny względem $\varphi$ nad $M$,
jeżeli istnieje formuła $d_{p,\varphi}(\bar y)$
z parametrami w $M$ taka, że dla każdej
krotki $\bar b$ z $M$ odpowiedniej długości
\[
\varphi(\bar x,\bar b)\in p
\quad\Longleftrightarrow\quad
\mathfrak M\models d_{p,\varphi}(\bar b).
\]
Typ nazywamy definiowalnym nad $M$, jeżeli
jest definiowalny względem każdej takiej formuły.
\end{definicja}

Definicja nie mówi, że sam typ jest izolowany
przez pojedynczą formułę w zmiennych $\bar x$.
Mówi, że dla ustalonej formuły zbiór parametrów,
dla których typ wybiera odpowiedź dodatnią,
jest definiowalny w modelu.

\begin{twierdzenie}[Definiowalność przy braku własności porządku]
\label{thm:mod-stable-definability}
Jeżeli żadna formuła języka $\mathcal L$
nie ma własności porządku w $T$, to każdy typ
zupełny skończonej krotki nad modelem $T$
jest definiowalny nad tym modelem.
\end{twierdzenie}
\begin{proof}
Ustalmy $\mathfrak M\models T$, $p\in S_n(M)$
oraz $\varphi(\bar x,\bar y)$.
Jeśli $\bar y$ jest pusta, wystarczy tautologia
lub sprzeczność, zależnie od tego, czy
$\varphi\in p$. Zakładamy dalej $|\bar y|\geq1$.
Dla krotek parametrów $\bar b$ połóżmy
$\epsilon(\bar b)=1$ wtedy i tylko wtedy,
gdy $\varphi(\bar x,\bar b)\in p$,
a w przeciwnym razie $\epsilon(\bar b)=0$.

Przypuśćmy, że nie istnieje skończona rodzina
krotek $\bar a_0,\ldots,\bar a_{k-1}$ z $M^n$,
dla której wektor wartości logicznych formuł
$\varphi(\bar a_i,\bar b)$ wyznacza
$\epsilon(\bar b)$. Zbudujemy indukcyjnie
krotki $\bar a_i,\bar b_i,\bar c_i$ w $M$.

Mając $\bar a_0,\ldots,\bar a_{j-1}$,
wybieramy krotki $\bar b_j,\bar c_j$
o identycznych wartościach logicznych
$\varphi(\bar a_i,\bar b_j)$ i
$\varphi(\bar a_i,\bar c_j)$ dla $i<j$,
lecz z $\epsilon(\bar b_j)=1$ oraz
$\epsilon(\bar c_j)=0$.
Jest to możliwe z przyjętego założenia,
również dla pustej rodziny prób.
Następnie wybieramy $\bar a_j\in M^n$
realizującą skończony fragment typu $p$:
\[
\bigwedge_{i=0}^j
\bigl(\varphi(\bar x,\bar b_i)\land
\neg\varphi(\bar x,\bar c_i)\bigr).
\]
Istnienie takiej krotki w $M$ wynika z
twierdzenia \ref{thm:mod-type-elementary-extension}:
każdy skończony fragment typu nad modelem
jest w tym modelu realizowany.

Dla formuły
\[
\psi(\bar x;\bar y,\bar z)=
(\varphi(\bar x,\bar y)
\leftrightarrow\varphi(\bar x,\bar z))
\]
otrzymujemy
\[
\mathfrak M\models\psi(\bar a_i;\bar b_j,\bar c_j)
\quad\Longleftrightarrow\quad i<j.
\]
Istotnie, dla $i<j$ wybór $\bar b_j,\bar c_j$
zapewnia zgodność wartości logicznych,
a dla $i\geq j$ wybór $\bar a_i$ zapewnia
ich przeciwne wartości.
Formuła $\psi$ ma więc własność porządku,
co jest sprzeczne z założeniem.

Istnieje zatem skończona rodzina prób
$\bar a_0,\ldots,\bar a_{k-1}$ wyznaczająca
$\epsilon(\bar b)$.
Niech $E\subseteq\{0,1\}^k$ będzie zbiorem
wektorów realizowanych przez krotki $\bar b$
z $\epsilon(\bar b)=1$. Definiujemy
\[
d_{p,\varphi}(\bar y)=
\bigvee_{\eta\in E}
\bigwedge_{i<k}\varphi(\bar a_i,\bar y)^{\eta(i)},
\]
gdzie $\theta^1=\theta$, $\theta^0=\neg\theta$.
Pusta alternatywa oznacza sprzeczność,
a pusta koniunkcja tautologię.
Z wyboru prób i definicji $E$ formuła ta
jest prawdziwa dokładnie dla parametrów
$\bar b$ z $\epsilon(\bar b)=1$.
Otrzymaliśmy wymaganą definiowalność.
\end{proof}

\begin{lemat}[Oszacowanie liczby typów definiowalnych]
\label{lem:mod-definability-count}
Niech $\tau=\max(\aleph_0,|\mathcal L|)$.
Jeżeli wszystkie typy jednozmienne nad modelem
$\mathfrak M\models T$ są definiowalne nad $M$,
a $|M|\leq\kappa$ dla $\kappa\geq\tau$, to
\[
|S_1(M)|\leq\kappa^\tau.
\]
\end{lemat}
\begin{proof}
Jest co najwyżej $\tau$ formuł języka
pierwotnego z wyróżnioną zmienną $x$
i skończoną krotką zmiennych parametrów.
Dla każdej z nich istnieje co najwyżej
$\kappa$ możliwych formuł definiujących
z parametrami z $M$, ponieważ język
$\mathcal L(M)$ ma co najwyżej $\kappa$
symboli, a formuły są skończonymi ciągami.

Dla każdego typu wybieramy rodzinę formuł
definiujących, po jednej dla każdej formuły
języka pierwotnego. Możemy ustalić dobre
uporządkowanie formuł i wybierać pierwszą
dopuszczalną definicję.
Rodzina wyznacza cały typ, gdyż każda
formuła z parametrami powstaje przez
podstawienie skończonej krotki parametrów
do formuły języka pierwotnego.
Liczba rodzin jest ograniczona przez
$\kappa^\tau$, co daje tezę.
\end{proof}

\begin{twierdzenie}[Kryterium stabilności]
\label{thm:mod-stability-iff-no-op}
Dla teorii zupełnej $T$ z modelami nieskończonymi
następujące warunki są równoważne:
\begin{enumerate}[label=(\roman*)]
\item $T$ jest stabilna;
\item żadna formuła nie ma własności porządku w $T$.
\end{enumerate}
Ponadto jeśli (ii) zachodzi i
$\tau=\max(\aleph_0,|\mathcal L|)$, to $T$
jest $\kappa$-stabilna dla każdej mocy
$\kappa\geq\tau$ spełniającej $\kappa^\tau=\kappa$.
\end{twierdzenie}
\begin{proof}
Implikacja (i)$\Rightarrow$(ii) wynika z
twierdzenia \ref{thm:mod-op-unstable}.
Załóżmy (ii) i wybierzmy moc $\kappa$
o podanych własnościach.
Dla dowolnego modelu $\mathfrak D\models T$
i $A\subseteq D$ o $|A|\leq\kappa$ stosujemy
dolne twierdzenie Löwenheima--Skolema
\ref{thm:downward-ls}. Otrzymujemy model
$\mathfrak M\preccurlyeq\mathfrak D$
zawierający $A$, o $|M|\leq\kappa$.
Każdy typ nad $A$ rozszerza się do typu
nad $M$: wybieramy jego realizację
w rozszerzeniu elementarnym $\mathfrak M$
na mocy twierdzenia
\ref{thm:mod-type-elementary-extension},
a następnie bierzemy typ tej realizacji nad $M$.
Wolno użyć $\mathfrak M$ zamiast $\mathfrak D$,
gdyż elementarność zapewnia równość teorii
w języku z nazwami elementów $A$.
Zatem odwzorowanie ograniczenia
$S_1(M)\to S_1(A)$ jest surjekcją.
Twierdzenie \ref{thm:mod-stable-definability}
daje definiowalność wszystkich typów nad $M$,
a lemat \ref{lem:mod-definability-count} daje
\[
|S_1(A)|\leq |S_1(M)|\leq\kappa^\tau=\kappa.
\]
Jest to $\kappa$-stabilność. Aby otrzymać
co najmniej jedną odpowiednią moc, bierzemy
$\kappa=2^\tau$, dla której
$(2^\tau)^\tau=2^{\tau\cdot\tau}=2^\tau$
z arytmetyki mocy nieskończonych.
\end{proof}

\begin{wniosek}[Niestabilność teorii porządku]
\label{cor:mod-dlo-unstable}
Teoria gęstych porządków liniowych bez końców
jest $\aleph_0$-kategoryczna, ale niestabilna
w każdej mocy nieskończonej.
\end{wniosek}
\begin{proof}
Kategoryczność wykazano w przykładzie
\ref{ex:mod-rn-dlo}. Formuła $x<y$ ma własność
porządku, jak pokazano w przykładzie
\ref{ex:mod-dlo-not-omega-stable}.
Stosujemy twierdzenie \ref{thm:mod-op-unstable}.
\end{proof}

Warunek $\kappa^\tau=\kappa$ jest warunkiem
wystarczającym, a nie koniecznym dla stabilności
w konkretnej mocy. Na przykład teoria
nieskończonych zbiorów z równością
(przykład \ref{ex:mod-stability-equality})
jest stabilna we wszystkich mocach nieskończonych,
także dla $\kappa=\aleph_0$, dla której
$\kappa^{\aleph_0}>\kappa$.
Dokładne spektrum stabilności wymaga
dalszych narzędzi.
